\chapter{Graphes, chemins et matrices d’adjacence}
\section*{Exercice 34.1 -- Degrés}
[Base] Un graphe simple non orienté a six sommets de degrés $1,2,2,3,3,3$. Combien possède-t-il d'arêtes ?

\section*{Exercice 34.2 -- Adjacence}
[Méthode] Écrire la matrice du chemin non orienté $A-B-C$, dans cet ordre. Calculer le nombre de marches de longueur 2 de $A$ à $A$ et de $A$ à $C$.

\section*{Exercice 34.3 -- Connexité}
[Raisonnement] Un graphe simple possède cinq sommets et deux composantes connexes de tailles 2 et 3. Quel est le nombre maximal d'arêtes ?

\section*{Exercice 34.4 -- Sommets de degré impair}
[Preuve] Montrer qu'un graphe fini non orienté possède un nombre pair de sommets de degré impair.

\section*{Exercice 34.5 -- Chemins et puissances}
[Preuve] Expliquer pourquoi $(M^{n+1})_{ij}=\sum_k(M^n)_{ik}M_{kj}$ compte les marches de longueur $n+1$.

\section*{Exercice 34.6 -- Orientation}
[Application] Dans le cycle orienté $A\to B\to C\to A$, combien de marches de longueur 2026 vont de $A$ à $B$ ?
